% Partie : incomplétude
% Chapitre : représentabilité dans Q
% Section : représentabilité des fonctions de C

\documentclass[../../../include/open-logic-section]{subfiles}

\begin{document}

\olfileid[fr]{inc}{req}{cre}
\olsection{Les fonctions de $C$ sont représentables dans $\Th{Q}$}

Nous devons montrer que toute fonction de $C$ est représentable
dans $\Th{Q}$. Il s'agit, en définitive, d'associer à chaque
fonction d'arité $k$, $f(x_0,\dots,x_{k-1})$, appartenant à $C$,
une formule $!A_f(x_0,\dots,x_{k-1},y)$ qui la représente.

\begin{lem}
Pour tous entiers naturels $n$ et $m$, si $n\neq m$, alors
$Q\Proves\num n\neq\num m$.
\end{lem}

\begin{proof}
Montrons par induction sur $n$ que, pour tout $m$, si
$n\neq m$, alors $Q\Proves\eq/[\num n][\num m]$.

Dans le cas de base, $n=0$. Si $m$ est différent de $0$,
alors $m=k+1$ pour un certain entier naturel $k$.
Nous disposons de l'axiome $\lforall[x][\eq/[0][x']]$.
Un axiome de quantification permet d'y remplacer $x$ par
$\num k$ et d'en déduire $\eq/[0][\num k']$.
Or $\num k'$ est précisément $\num m$.

Pour l'étape d'induction, supposons la propriété établie
pour $n$ et considérons $n+1$. Soit $m$ un entier naturel
quelconque. Deux cas sont possibles : soit $m=0$, soit
$m=k+1$ pour un certain $k$. Le premier cas se traite comme
ci-dessus. Dans le second, supposons que $n+1\neq k+1$.
Alors $n\neq k$. L'hypothèse d'induction pour $n$ donne
$\Th{Q}\Proves\eq/[\num n][\num k]$.
Nous disposons de l'axiome
$\lforall[x][\lforall[y][\eq[x'][y']\lif\eq[x][y]]]$.
Un axiome de quantification donne
$\eq[\num n'][\num k']\lif\eq[\num n][\num k]$.
La logique propositionnelle permet alors de déduire, dans
$\Th{Q}$,
$\eq/[\num n][\num k]\lif\eq/[\num n'][\num k']$.
Par modus ponens, nous obtenons
$\eq/[\num n'][\num k']$, ce qui est le résultat voulu,
puisque $\num k'$ est $\num m$.
\end{proof}

Le lemme affirme, au fond, assez peu de chose :
$\Th{Q}$ peut démontrer que des numéraux distincts désignent
des objets distincts. Par exemple, $\Th{Q}$ démontre
$0''\neq 0'''$. Il faut néanmoins procéder avec soin pour
l'établir en général. Remarquons aussi que l'induction
employée ici est une induction \emph{extérieure} à
$\Th{Q}$.

Nous pourrons représenter la fonction nulle, le successeur,
l'addition, la multiplication, la fonction caractéristique
de l'égalité et les projections. Dans chaque cas, la formule
représentante est très simple. Par exemple, la fonction
nulle est représentée par la formule
\[
y = 0,
\]
le successeur par la formule
\[
x_0' = y,
\]
et l'addition par la formule
\[
x_0 + x_1 = y.
\]
Le travail consiste à montrer que $\Th{Q}$ peut démontrer
les énoncés requis. Ainsi, affirmer que la formule précédente
représente l'addition revient à montrer que, pour tous
entiers naturels $m$ et $n$, $\Th{Q}$ démontre
\[
\eq[\num n + \num m][\num {n+m}]
\]
et
\[
\lforall[y][(\eq[(\num n + \num m)][y] \lif \eq[y][\num {n+m}])].
\]

Qu'en est-il de la composition ? Supposons que $h$ soit
définie par
\[
h(x_0,\dots,x_{l-1}) = f(g_0(x_0,\dots,x_{l-1}), \dots,
g_{k-1}(x_0,\dots,x_{l-1})).
\]
Nous avons déjà trouvé des formules
$!A_f, !A_{g_0}, \dots, !A_{g_{k-1}}$ représentant
respectivement les fonctions $f,g_0,\dots,g_{k-1}$.
Nous pouvons alors définir une formule $!A_h$ représentant
$h$ en posant $!A_h(x_0,\dots,x_{l-1},y)$ égale à
\begin{multline*}
  \lexists[z_0\dots][\lexists[z_{k-1}][(!A_{g_0}(x_0,\dots,x_{l-1},z_0) \land
  \dots \land !A_{g_{k-1}}(x_0,\dots,x_{l-1},z_{k-1}) \land]]\\
  !A_f(z_0,\dots,z_{k-1},y)).
\end{multline*}

Considérons enfin la recherche non bornée.
Supposons que $g(x,\vec z)$ soit régulière et représentable
dans $\Th{Q}$, par exemple par la formule
$!A_g(x,\vec z,y)$. Définissons $f$ par
$f(\vec z)=\mu x\;g(x,\vec z)$.
Nous cherchons une formule $!A_f(\vec z,y)$ qui représente
$f$. Le choix suivant est naturel :
\[
!A_f(\vec z,y) \ident !A_g(y,\vec z,0) \land \lforall[w][(w < y
\lif \lnot !A_g(w,\vec z,0))].
\]

On peut montrer que cette formule convient à l'aide de
quelques lemmes supplémentaires, par exemple :

\begin{lem}
Pour toute variable $x$ et tout entier naturel $n$,
$\Th{Q}$ démontre $\eq[(x'+\num n)][(x+\num n)']$.
\end{lem}

Cet énoncé pour chaque numéral fixé est plus faible que
l'énoncé universel
$\lforall[x][\lforall[y][(x'+y)=(x+y)']]$,
que $\Th{Q}$ ne démontre pas.

\begin{proof}
La démonstration se fait, comme d'habitude, par induction
sur $n$. Dans le cas de base, $n=0$, nous devons montrer
que $\Th{Q}$ démontre $\eq[(x'+0)][(x+0)']$.
Or nous avons :
\begin{eqnarray*}
& & \eq[(x' + 0)][x'] \quad \text{par l'axiome 4} \\
& & \eq[(x + 0)][x] \quad \text{par l'axiome 4} \\
& & \eq[(x+0)'][x'] \quad \text{par la ligne 2} \\
& & \eq[(x' + 0)][(x + 0)'] \quad \text{par les lignes 1 et 3}
\end{eqnarray*}
Pour l'étape d'induction, supposons que nous ayons démontré
$\eq[(x'+\num n)][(x+\num n)']$ dans $\Th{Q}$.
Puisque $\num{n+1}$ est $\num n'$, il nous faut montrer
que $\Th{Q}$ démontre
$\eq[(x'+\num n')][(x+\num n')']$.
La chaîne d'égalités suivante est démontrable dans $\Th{Q}$ :
\begin{eqnarray*}
(x' + \num n') & \eq & (x' + \num n)' \quad \text{axiome 5} \\
& \eq & (x + \num n')' \quad \text{hypothèse d'induction et axiome 5}
\end{eqnarray*}
La dernière ligne utilise aussi l'axiome 5 pour réécrire
le successeur de la somme.\footnote{Précision de l'édition
française : la source ne mentionne que l'hypothèse
d'induction pour cette ligne ; l'axiome 5 est également
nécessaire.}
\end{proof}

\begin{lem}
\begin{enumerate}
\item $\Th{Q}$ démontre $\lnot(x<\num 0)$.
\item Pour tout entier naturel $n$, $\Th{Q}$ démontre
\[
x < \num {n+1} \lif (\eq[x][\num 0] \lor \dots \lor \eq[x][\num n]).
\]
\end{enumerate}
\end{lem}

\begin{proof}
Traitons le point 1 et une partie du point 2 de façon
informelle : nous donnons seulement des indications pour
construire la !!{derivation} formelle.

Pour le point 1, la définition de $<$ nous ramène à
démontrer $\lnot\lexists[y][\eq[(y'+x)][0]]$ dans
$\Th{Q}$. D'après les axiomes et les règles de la logique
du premier ordre, cela équivaut à
$\lforall[y][\eq/[(y'+x)][0]]$.
Voici l'idée : supposons que $\eq[(y'+x)][0]$.
Si $x$ vaut 0, nous avons $\eq[(y'+0)][0]$.
Mais l'axiome 4 de $\Th{Q}$ donne $\eq[(y'+0)][y']$,
tandis que l'axiome 2 donne $\eq/[y'][0]$ :
nous obtenons une contradiction.
Si $x$ est différent de $0$, l'axiome 3 fournit un $z$
tel que $\eq[x][z']$. Nous avons alors
$\eq[(y'+z')][0]$ ; l'axiome 5 donne
$\eq[(y'+z)'][0]$, ce qui contredit de nouveau
l'axiome 2. Les deux cas conduisent donc à une contradiction,
et nous pouvons conclure
$\lforall[y][\eq/[(y'+x)][0]]$.\footnote{Correction de
l'édition française : la source place cette conclusion
universelle après le premier cas seulement. Elle est ici
placée après l'examen des deux cas.}

Pour le point 2, procédons par induction sur $n$.
Considérons le cas de base, $n=0$.
Il faut montrer que $x<\num 1\lif\eq[x][\num 0]$.
Supposons que $x<\num 1$. L'axiome définissant $<$ donne
$\lexists[y][\eq[(y'+x)][0']]$.
Soit $y$ possédant cette propriété ; nous avons donc
$\eq[y'+x][0']$.

Nous devons montrer que $\eq[x][0]$.
Si $x$ est différent de 0, l'axiome 3 indique qu'il est
égal à $z'$ pour un certain $z$. Nous avons alors
$\eq[(y'+z')][0']$.
L'axiome 5 de $\Th{Q}$ donne
$\eq[(y'+z)'][0']$, puis l'axiome 1 donne
$\eq[(y'+z)][0]$. Par définition, cela signifie
$z<0$, contrairement au point 1.
\end{proof}

\begin{explain}
Nous avons caractérisé les fonctions calculables comme
les fonctions représentables dans $\Th{Q}$.
L'argument a, en réalité, une portée plus générale.
D'après la définition de la représentabilité, toute théorie
qui étend $\Th{Q}$, ou dans laquelle on peut interpréter
$\Th{Q}$, peut représenter les fonctions calculables.
Réciproquement, si une théorie est cohérente, distingue
démontrablement les numéraux distincts et dispose d'un système
de !!{derivation}s dont les preuves sont effectivement
vérifiables, toute fonction qu'elle représente est
calculable.\footnote{Qualification de l'édition française :
la source omet ici la cohérence de la théorie et la distinction
démontrable des numéraux. Ces conditions garantissent que
la recherche effective de preuves retrouve la valeur correcte.
Une théorie incohérente satisfait trivialement les conditions
de représentabilité pour toute fonction, même non calculable.
Les caractérisations qui suivent supposent ces conditions.}
Ainsi, les fonctions calculables peuvent aussi être
caractérisées comme celles qui sont représentables dans
l'arithmétique de Peano, ou même dans la théorie des
ensembles de Zermelo--Fraenkel.
Comme Gödel l'a remarqué, ce résultat peut surprendre.
Nous verrons que les questions de démontrabilité dépendent
étroitement de la théorie considérée : en gros, des axiomes
plus forts permettent de démontrer davantage.
Pourtant, dans une large classe de théories axiomatiques
satisfaisant les conditions précédentes, les fonctions
représentables sont exactement les fonctions calculables.
\end{explain}

\end{document}
